\section{The Causality Game — Formal Framework}
\label{sec:formaldefinition}
\game is not a metaphor but a formal interactive system for causal discovery and reasoning grounded in structural causal models (SCMs). 
It represents a controlled, turn-based environment in which an agent performs interventions, observes results, and progressively uncovers causal relationships among a finite set of variables.
The following formulation describes the game as a mathematical object, providing the foundation upon which its software implementation is built.

\subsection{Definition and Core Components}
A single game instance is defined by the tuple
\[
    \langle 
        \scm,\; \latentx \sqcup \measuredx \sqcup \controllablex,\;
        \initstate,\; \phi,\;
        \goal,\; \resultmetric,\; \behaviormetric
    \rangle ,
\]
where:
\begin{itemize}
    \item $\scm$ is the underlying structural causal model that governs the world’s data-generating process.
    \item $\latentx, \measuredx, \controllablex$ partition $\varset$, determining which variables are hidden, observable, or intervenable.
    \item $\initstate \in \statespace$ represents the starting data context, typically empty or initialized with baseline observations under the null intervention $\bot$.
    \item $\phi : \runs \to \{0,1\}$ is a stopping function that enforces runtime limits such as sample count, memory, or time budget.
    \item $\goal$ defines the mission — the causal question to be solved.
    \item $\resultmetric : \outputspace \to \mathbb{R}$ scores the correctness of the final answer.
    \item $\behaviormetric : \runs \to \mathbb{R}$ measures the efficiency of the agent’s experimentation behavior.
\end{itemize}

The agent does not directly observe $\scm$ or any $\latentx$; instead, it interacts through interventions on $\controllablex$ and receives measured outcomes from $\measuredx \sqcup \controllablex$ .

\subsection{Extended Structural Causal Model}
Formally, the SCM as the foundation of the game, is defined as
\[
    \scm = \langle \varset, \{f_i\}, \{\eta_i\} \rangle,
\]
where $\varset$ is the set of variables, $f_i$ are deterministic structural functions, and $\eta_i$ are exogenous noise terms.
In the context of \game, each variable also has a declared \emph{access mode}:
\[
    \latentx \sqcup \measuredx \sqcup \controllablex = \varset.
\]
This configuration explicitly controls which nodes the agent may observe or intervene upon.  
The SCM thus defines both the ground-truth causal structure and the constraints under which agents can interact with it, reflecting real-world experimental limitations.

\subsection{Environment and Experiment Execution}
The environment is a deterministic causal oracle induced by the extended SCM. It does not maintain any mutable internal ``state.’’ Instead, it responds to each agent action by \emph{functionally generating new data} and returning it to the agent. All memory of past interactions is externalized into the \textbf{Transcript}, which is maintained by the \game orchestrator, not by the environment itself.

\paragraph{Interventions.}
An intervention is a partial assignment over controllable variables:
\[
    \intervention : X_i \in \controllablex \;\mapsto\; v \in domain(X_i) \;\cup\; \{\none\},
\qquad
\intervention \in \interventions,
\]
where $\intervention(X_i) = \none$ means that $X_i$ is left to follow its natural SCM equation. The empty intervention $\bot$ corresponds to purely observational sampling.

\paragraph{Experiments.}
Each \emph{environment action} is a finite experiment request:
\[
    e \;:\; \interventions \;\to\; \mathbb{N},
\]
meaning that for each intervention $\intervention$ with $e(\intervention) = n > 0$, the environment returns $n$ new i.i.d.\ samples generated under that interventional condition. There is no deferred scheduling: each experiment is executed immediately in the same round it is requested.

\paragraph{Interaction Protocol (inform $\rightarrow$ act).}
The agent never receives a raw environment state. Instead, the interaction in each round proceeds in two stages:
\begin{enumerate}
    \item \textbf{inform:} The agent is given the newly generated samples from its previous experiment, restricted to $\measuredx \cup \controllablex$, together with a structured feedback object (e.g.\ partial scoring about how well its current hypothesis aligns with the mission).
    \item \textbf{act:} The agent is then given two structured DTOs:
    \begin{itemize}
        \item a \texttt{RoundInformation} object summarizing the current budget status, mission description, etc.,
        \item an \texttt{AvailableActions} object enumerating all legal next experiment configurations plus the special terminal option \emph{stop with answer}.
    \end{itemize}
    It must then choose either another experiment $e \in \experiments$ or stop and submit a final answer.
\end{enumerate}


\subsection{Environment Determinism and Reproducibility}
The agent’s experience is not merely stochastic — it is \emph{seed-deterministic}. The initialization seed fixes the entire sample-generation randomness of the environment. Thus, two independent agents that issue the same sequence of actions will observe \emph{bitwise identical outcomes}. This ensures \textbf{perfect fairness and reproducibility} across agents and research groups — a critical requirement for benchmark integrity.

\subsection{Start and Stop Conditions}
The temporal extent of a run is controlled by two mechanisms:

\begin{itemize}
    \item an explicit \textbf{initial state} $\initstate \in \statespace$, which determines what data (if any) the agent sees before taking its first action,
    \item an external \textbf{stopping function} $\phi : \runs \to \{0,1\}$, which may forcibly terminate the game if resource constraints are exceeded (e.g.\ too many interventions or samples have been requested).
\end{itemize}

Crucially, the agent is always free to stop \emph{voluntarily} at any round by issuing a \emph{stop with answer} action, returning some $y \in \outputspace$. If $\phi$ triggers first, the game ends immediately \emph{without} returning an answer — resulting in a zero or penalty score under $\resultmetric$.

\paragraph{Scientific motivation.}
This design mirrors real experimental science:  
if a researcher exhausts funding before returning a final conclusion, there is \emph{no result to evaluate}.  
Thus, \game formalizes causal inference not only as a problem of correctness — but of disciplined experimental planning under pressure.


This reflects real-world constraints: experimenters may choose to conclude early, but if they exceed allowable resources, the process is terminated before any formal conclusion can be submitted.

\subsection{Capabilities and Limits}
Within each run, the agent:
\begin{itemize}
    \item can perform any legal intervention $\intervention \in \interventions$ on $\controllablex$,
    \item can request $n$ samples per intervention via experiments,
    \item receives only observable outcomes $\measuredx$,
    \item and may stop at any point with an answer submission.
\end{itemize}
It cannot:
\begin{itemize}
    \item access or view latent variables $\latentx$,
    \item modify the true $\scm$ or observe its internal structure,
    \item or bypass the stopping rule $\phi$ once triggered.
\end{itemize}
This strict separation between what the agent \emph{can} and \emph{cannot} know enforces an epistemically valid experimental process.

\subsection{Objective}
A \emph{player agent} interacts with the game instance by selecting actions in discrete turns to generate a run
\[
    r = (\initstate, a_0, s_1, a_1, \ldots, a_{T-1}, s_T),
    \quad r \in \runs.
\]
Its goal is to produce a final answer $y \in \outputspace$, by issuing a \emph{stop with answer} action before termination, such that $\resultmetric(y)$ is maximized while incurring minimal cost under $\behaviormetric(r)$.

Formally, a possible multi-objective evaluation to maximize and rank agents:
\[
    \text{score}(r) \;=\; 
    \big( \resultmetric(y),\;\; -\behaviormetric(r) \big),
\quad \text{Pareto-optimality is well-defined.}
\]
However. No scalarization is imposed at the framework level — users may rank agents by convex combination, lexicographic ordering, or full Pareto frontier analysis.

\subsection{Narrative Summary}
\game thus formalizes the process of causal discovery as an interactive and bounded experimentation problem.  
At each round, the agent receives structured contextual information, selects from permissible actions, and shapes the data it will later use to infer causal structure.  
Through this formulation, causal learning becomes a reproducible, quantitative, and comparable process — bridging the conceptual gap between theoretical causal reasoning and experimental design.